
\subsection{Combining Fixture Scopes}

Using two fixtures, you can combine the benefits of caching and isolation:

\begin{minted}[autogobble, escapeinside=||,highlightlines={1,8}]{python}
@pytest.fixture(scope="module")
def rpn_instance() -> RPNCalc:
    time.sleep(2)
    return RPNCalc(Config())

@pytest.fixture
def rpn(rpn_instance: RPNCalc):
    rpn_instance.stack.clear()
    return rpn_instance

def test_a(rpn: RPNCalc):
    rpn.stack.append(42)
    assert rpn.stack == [42]

def test_b(rpn: RPNCalc):
    assert not rpn.stack
\end{minted}

\mono{rpn_instance} uses caching and does some sort of slow operation
(initializing database, starting external process, initializing a device, etc.),
but it's not used directly in tests. Instead, \mono{rpn} is used, which in some
way resets the state of the underlying object before every test.

\begin{center}
    \begin{tikzpicture}
        \node [fixture, minimum width=3cm] (instance) {\mono{rpn_instance}};
        \node [fixtureval, below=5mm of instance] (rpnval) {\scriptsize \phantom{\mono{rpn()}}};
        \node [above=0.6mm of rpnval.north east, gray, xshift=0.2mm] {\scriptsize (2s)};
        \draw [-] (instance) -- (rpnval);

        \node [fixture, minimum width=2cm, below=5mm of rpnval, xshift=3em] (rpna) {\mono{rpn}};
        \node [fixture, minimum width=2cm, below=5mm of rpnval, xshift=-3em] (rpnb) {\mono{rpn}};
        \node [test, minimum width=2cm, below=5mm of rpna] (testa) {\mono{test_a}};
        \node [test, minimum width=2cm, below=5mm of rpnb] (testb) {\mono{test_b}};
        \draw [->] (rpna) -- (rpnval);
        \draw [->] (rpnb) -- (rpnval);
        \draw [->] (testa) -- (rpna);
        \draw [->] (testb) -- (rpnb);
    \end{tikzpicture}
\end{center}
\pagebreak
